\documentclass{note}
\title{Glass bead game}
\date{2026-04-29}
\modified{2026-04-29}
\begin{document}

\emph{Glass bead game} (German: \emph{Das Glasperlenspiel}) is a fictional game from Herman Hesse's book with the same title.
The books describes the game as a mix of mathematics and music played with beads, but doesn't give its rules.

\blockquote{
  Under the shifting hegemony of now this, now that science or art,
  the Game of games had developed into a kind of universal language through which the players could express values and set these in relation to one another.
  Throughout its history the Game was closely allied with music,
  and usually proceeded according to musical or mathematical rules.
  One theme, two themes, or three themes were stated, elaborated, varied, and underwent a development quite similar to that of the theme in a Bach fugue or a concerto movement.
  A Game, for example, might start from a given astronomical configuration, or from the actual theme of a Bach fugue, or from a sentence out of Leibniz or the Upanishads,
  and from this theme, depending on the intentions and talents of the player,
  it could either further explore and elaborate the initial motif or else enrich its expressiveness by allusions to kindred concepts.
  Beginners learned how to establish parallels, by means of the Game's symbols, between a piece of classical music and the formula for some law of nature.
  Experts and Masters of the Game freely wove the initial theme into unlimited combinations.
}{Herman Hesse, \cite{The Glass Bead Game}}

Articles tagged \ref{glasperlenspiel}{#glasperlenspiel} juxtapose ideas from different domains.

\end{document}
